% ============================================================================
% Fichier  : partie3/P03_C15_rapport.tex
% Titre    : Le Rapport et le Suivi des Actions
% Auteur   : MINKA MI NGUIDJOÏ Thierry Emmanuel
% Version  : 1.0 -- novembre 2026
% ============================================================================
% Sources : ISO 19011:2018 (6.5 et 6.7) ; ISO/IEC 27001:2022 (10.2).
% Rapport d'audit d'octobre 2026 de LiZbiZ : cas fictif.
% ============================================================================

\chapter{Le Rapport et le Suivi des Actions}
\label{chap:rapport}
\chaptermark{Rapport et suivi}

\epigraph{%
    « Un audit sans suivi\\
    n'est qu'un inventaire des regrets. »%
}{Principe du cours}

\niveauChapitre{Intermédiaire}{Chapitres~\ref{chap:audit}
à~\ref{chap:auditer-risques}.}

\begin{boxobjectifs}
\begin{itemize}
    \item Structurer un rapport d'audit pour ses différents lecteurs.
    \item Hiérarchiser les constats et formuler des recommandations
          vérifiables.
    \item Recueillir et traiter la réponse de l'audité.
    \item Conduire une action corrective selon l'ISO/IEC~27001 (10.2).
    \item Suivre la clôture des actions et en vérifier l'efficacité.
\end{itemize}
\end{boxobjectifs}

\bigskip

% ============================================================================
\section*{Question directrice : les constats de 2023 avaient-ils été suivis~?}
% ============================================================================

En relisant le rapport d'audit de 2023, Mme~NKENG découvre une
recommandation~: « étendre le second facteur aux accès des tiers ». Elle
avait été acceptée, puis oubliée. Aucun suivi n'avait été organisé. Ce
chapitre montre comment écrire un rapport qui sera lu, et organiser un
suivi qui sera tenu.

\medskip
\textbf{Plan du chapitre.}
\begin{enumerate}
    \item La structure du rapport (section~\ref{sec:c15-structure}).
    \item Hiérarchiser et recommander (section~\ref{sec:c15-recommander}).
    \item La réponse de l'audité (section~\ref{sec:c15-reponse}).
    \item L'action corrective (section~\ref{sec:c15-corrective}).
    \item Suivre et clore (section~\ref{sec:c15-suivi}).
\end{enumerate}

\separateur

% ============================================================================
\section{La Structure du Rapport}
\label{sec:c15-structure}
% ============================================================================

\begin{tableaucours}{L{3.6cm}L{8.6cm}}{%
    \tableauHeader{Partie} &
    \tableauHeader{Contenu}
}
Synthèse & Une page pour la direction~: conclusion d'ensemble, trois à
cinq constats majeurs, décisions demandées \\
Cadre & Objectifs, périmètre, critères, période, équipe, limites \\
Constats & Chaque constat avec critère, fait, preuve, cause, conséquence,
classement (Chapitre~\ref{chap:conduire}) \\
Recommandations & Une ou plusieurs par constat, vérifiables \\
Réponse de l'audité & Accord ou désaccord, plan d'action, responsable,
échéance \\
Annexes & Liste des preuves, échantillons, personnes rencontrées \\
\end{tableaucours}
\vspace{-0.8em}
\begin{center}{\small\itshape Tableau~15.1. Structure d'un rapport d'audit}\end{center}

\begin{boxalert}
\textbf{Violation typique.} Taire les limites de l'audit.

\smallskip
\textbf{Règle~:} le rapport dit ce qui n'a pas été examiné et pourquoi
(accès refusé, temps, périmètre). Un lecteur qui croit tout couvert prend
les silences pour des conformités.
\end{boxalert}

% ============================================================================
\section{Hiérarchiser et Recommander}
\label{sec:c15-recommander}
% ============================================================================

\begin{tableaucours}{L{2.4cm}L{5.6cm}L{4.2cm}}{%
    \tableauHeader{Priorité} &
    \tableauHeader{Critère} &
    \tableauHeader{Délai de traitement}
}
1 & Non-conformité majeure~; risque intolérable non réduit & Un mois \\
2 & Non-conformité mineure~; risque à traiter & Trois mois \\
3 & Opportunité d'amélioration & Selon le plan annuel \\
\end{tableaucours}
\vspace{-0.8em}
\begin{center}{\small\itshape Tableau~15.2. Hiérarchiser les constats (convention du cours)}\end{center}

Une recommandation \emph{vérifiable} dit ce qui doit être vrai, et non ce
qu'il faudrait « envisager »~: « au 30~novembre, aucun compte d'accès
distant ne se connecte sans second facteur, et l'indicateur est calculé
sur toutes les connexions » se vérifie~; « renforcer la sécurité des
accès » ne se vérifie pas.

\begin{boxscenario}{Opération LIZBIZ. La synthèse du rapport d'octobre}
\textbf{Conclusion d'ensemble.} Le dispositif de gestion des risques est
correctement conçu et fonctionne en partie~; son efficacité est
surestimée.

\medskip
\textbf{Constats majeurs.} A-03, compte d'accès distant sans second
facteur (priorité~1)~; A-05, indicateur calculé sur une population
incomplète (priorité~2)~; A-07, deux risques résiduels acceptés sans
trace (priorité~2)~; A-08, deux risques manquants au registre
(priorité~2).

\medskip
\textbf{Décision demandée.} Supprimer toute exception au second facteur
sans validation du comité des risques.
\end{boxscenario}

% ============================================================================
\section{La Réponse de l'Audité}
\label{sec:c15-reponse}
% ============================================================================

L'audité répond à chaque constat~: il l'accepte ou le conteste, et propose
une action, un responsable et une échéance. Un désaccord n'est pas effacé~:
il figure au rapport, avec les arguments des deux parties (principe de
présentation impartiale, Chapitre~\ref{chap:audit}).

\begin{tableaucours}{L{1.2cm}L{2.4cm}L{5.0cm}L{1.6cm}L{1.6cm}}{%
    \tableauHeader{Constat} &
    \tableauHeader{Réponse} &
    \tableauHeader{Action} &
    \tableauHeader{Resp.} &
    \tableauHeader{Échéance}
}
A-03 & Accepté & Second facteur sur « vpn\_secours », procédure de
dérogation & M.~KOUNOU & 30~nov. \\
A-05 & Accepté & Recalcul de l'indicateur sur toutes les connexions &
Mme~MOUAFO & 15~nov. \\
A-07 & Accepté & Formulaire d'acceptation signé & Mme~MOUAFO & 30~nov. \\
A-08 & Partiellement & Ajout du risque « équipement nomade »~; la
climatisation est jugée couverte par R-01 & Mme~MOUAFO & 31~déc. \\
\end{tableaucours}
\vspace{-0.8em}
\begin{center}{\small\itshape Tableau~15.3. Réponse de la direction au rapport d'octobre}\end{center}

% ============================================================================
\section{L'Action Corrective}
\label{sec:c15-corrective}
% ============================================================================

\begin{boxprocedure}
\textbf{Action corrective (d'après l'ISO/IEC 27001, 10.2)}
\begin{enumerate}
    \item \textbf{Réagir}~: maîtriser et corriger l'écart (activer le second
          facteur sur le compte).
    \item \textbf{Analyser la cause}~: pourquoi l'écart existe-t-il, et
          peut-il exister ailleurs~? (exceptions orales non tracées).
    \item \textbf{Agir sur la cause}~: procédure de dérogation écrite,
          limitée dans le temps, validée par le comité.
    \item \textbf{Vérifier l'efficacité}~: refaire l'essai, recalculer
          l'indicateur, rechercher d'autres exceptions.
    \item \textbf{Mettre à jour} le registre et, si besoin, le système de
          management.
\end{enumerate}
\end{boxprocedure}

\begin{boxdef}
\termecle{Correction} et \termecle{action corrective}~: la correction
supprime l'écart constaté~; l'action corrective supprime sa cause pour
éviter qu'il ne se reproduise.
\end{boxdef}

% ============================================================================
\section{Suivre et Clore}
\label{sec:c15-suivi}
% ============================================================================

Le suivi est tenu par l'audit (troisième ligne) dans un tableau des
recommandations, présenté au comité d'audit. Une action n'est
\emph{close} que lorsque l'auditeur en a vérifié l'efficacité, et non
lorsque l'audité la déclare faite.

\begin{boxobjet}{Constat A-03 · Suivi}
\textbf{Action~:} second facteur sur le compte « vpn\_secours »,
procédure de dérogation. \textbf{Déclarée réalisée~:} 22~novembre.
\textbf{Vérifiée~:} 3~décembre, par un essai de connexion (refusé sans
second facteur) et par l'export du journal (aucune connexion sans second
facteur depuis le 22~novembre). \textbf{Statut~:} close.
\end{boxobjet}

\begin{tableaucours}{L{6.0cm}R{3.0cm}}{%
    \tableauHeader{Indicateur de suivi} &
    \tableauHeader{Fin décembre}
}
Recommandations émises & 9 \\
Closes après vérification & 6 \\
En cours, dans les délais & 2 \\
En retard & 1 \\
\end{tableaucours}
\vspace{-0.8em}
\begin{center}{\small\itshape Tableau~15.4. Tableau de suivi des recommandations (audit d'octobre)}\end{center}

\begin{boxmetier}{Auditeur des systèmes d'information}
\textbf{Ne clore que ce qu'on a vérifié.} Mme~NKENG refuse de clore sur
une déclaration~: elle demande la preuve de mise en œuvre, puis la preuve
d'efficacité. La recommandation de 2023, déclarée « en cours » pendant
trois ans, ne l'aurait pas été sous ce régime.
\end{boxmetier}

\separateur

\begin{boxresume}
\textbf{Ce qu'il faut retenir de ce chapitre~:}
\begin{itemize}
    \item \textbf{Le rapport} s'adresse à plusieurs lecteurs~: synthèse,
          cadre et limites, constats, recommandations, réponse, annexes.
    \item \textbf{Les constats} se hiérarchisent~; les recommandations sont
          vérifiables.
    \item \textbf{La réponse de l'audité} figure au rapport, désaccords
          compris.
    \item \textbf{L'action corrective} traite la cause, pas seulement
          l'écart, et vérifie son efficacité.
    \item \textbf{Une action est close} quand l'auditeur a vérifié son
          efficacité.
\end{itemize}
\end{boxresume}

\begin{boxinfo}
\textbf{Points de contrôle de l'élève.} Je sais~: structurer un rapport~;
hiérarchiser des constats~; rédiger une recommandation vérifiable~;
distinguer correction et action corrective~; tenir un tableau de suivi.
\end{boxinfo}

\bigskip

\begin{boxexo}[title={\ding{45}\quad Exercice~15.1. Vérifiable \niveaubase}]
Réécrivez de manière vérifiable~: « améliorer la gestion des sauvegardes ».
\end{boxexo}

\begin{boxexo}[title={\ding{45}\quad Exercice~15.2. Priorité \niveaubase}]
Donnez la priorité des constats A-03, A-05 et A-08 selon le
Tableau~15.2, et le délai correspondant.
\end{boxexo}

\begin{boxexo}[title={\ding{45}\quad Exercice~15.3. Correction ou action corrective \niveaubase}]
Classez~: activer le second facteur sur le compte~; créer une procédure de
dérogation~; supprimer le compte.
\end{boxexo}

\begin{boxexo}[title={\ding{45}\quad Exercice~15.4. La synthèse \niveaumoyen}]
Réécrivez la synthèse du rapport d'octobre en dix lignes pour le conseil.
\end{boxexo}

\begin{boxexo}[title={\ding{45}\quad Exercice~15.5. Le désaccord \niveaumoyen}]
La direction conteste en partie A-08. Comment le rapport doit-il présenter
ce désaccord~? L'auditrice doit-elle céder~?
\end{boxexo}

\begin{boxexo}[title={\ding{45}\quad Exercice~15.6. Les limites \niveaumoyen}]
Rédigez la section « limites » du rapport d'octobre~: ce qui n'a pas été
examiné et pourquoi.
\end{boxexo}

\begin{boxexo}[title={\ding{45}\quad Exercice~15.7. Clore \niveaumoyen}]
M.~KOUNOU déclare l'action A-03 réalisée. Quelles preuves Mme~NKENG
exige-t-elle avant de la clore~?
\end{boxexo}

\begin{boxexo}[title={\ding{45}\quad Exercice~15.8. Le suivi \niveaumoyen}]
Calculez le taux de clôture vérifiée du Tableau~15.4. Que proposez-vous
pour l'action en retard~?
\end{boxexo}

\begin{boxexo}[title={\ding{45}\quad Exercice~15.9. La recommandation de 2023 \niveauexpert}]
Analysez pourquoi la recommandation de 2023 a été oubliée, et proposez un
dispositif de suivi qui l'aurait empêché.
\end{boxexo}

\begin{boxexo}[title={\ding{45}\quad Exercice~15.10. Un rapport complet \niveauexpert}]
Rédigez le rapport complet de l'audit « accès distants et prestataire »,
selon la structure du Tableau~15.1.
\end{boxexo}

\bigskip

\begin{boxinfo}
\textbf{Pour aller plus loin.}
\begin{itemize}
    \item ISO~19011, rapport et suivi~\cite{iso19011_2018}.
    \item ISO/IEC~27001, clause~10~\cite{iso27001_2022}.
\end{itemize}
\end{boxinfo}

\begin{boxinfo}
\textbf{Transition.} L'audit vérifie que les mesures tiennent. Quand
elles cèdent, il faut gérer l'incident. La Partie~IV commence par la
gestion des incidents (Chapitre~\ref{chap:incidents}).
\end{boxinfo}
